\section{Participant Eligibility, Session Schedule, and Protocol Fidelity}

\subsection{Eligibility and recruitment}
The intended population comprises school-age children with an independently documented ADHD diagnosis or another explicitly defined research eligibility pathway. Diagnostic status, medication regimen, relevant co-occurring conditions, and recruitment source will be documented using a prespecified protocol; an in-game score cannot establish diagnosis. Age range, inclusion criteria, exclusion criteria, and recruitment setting require approval before enrollment. The study will not presume that every child with ADHD can or should complete an AR task.

Participation requires caregiver permission and age-appropriate child assent. Refusal or withdrawal is permitted at any point without penalty. Exclusion from the AR substudy should be limited to inability to conduct the specific procedure safely or reliably after reasonable accessibility adjustments, rather than broad exclusion based on neurodevelopmental differences. Reasons for screening failure and module-specific nonqualification will be reported in aggregate.

\subsection{Screening and baseline}
Screening will verify the relevant eligibility criteria, caregiver and child understanding of the procedure, safety considerations, and compatibility of the family-owned Android phone. A brief supervised familiarization session will establish whether the child understands each cognitive rule independently of response modality. Device qualification will test inertial input, AR tracking, timing stability, and a safe operating area. Baseline assessment will precede randomization, and its primary score will not be used to selectively admit children based on anticipated treatment response.

\subsection{Proposed sequence of visits}
The study proceeds through a screening and familiarization visit, baseline transfer assessment, training sessions, an immediate post-training assessment, and a retention assessment. The precise number and spacing of sessions, maximum block duration, inter-block breaks, and follow-up interval are design parameters to be frozen following feasibility work. They are not presented as established doses.

Within a training visit, the sequence is a brief comfort check, device and environment check, standardized task instructions, practice trials, eligible challenge blocks, child-controlled breaks, and an end-of-session tolerability check. The controller records challenge, regulation, and recovery events under the same core definitions across contexts. The child may stop before the planned session endpoint. A missed or shortened session is recorded with its reason and is not replaced without a prespecified rule.

\subsection{Allocation and masking}
For the initial matched-context experiment, condition order will be randomized or counterbalanced within participant while respecting safety and washout rules. For the later transfer trial, randomization is between fixed-context and varied-context training schedules. Allocation generation, concealment, stratification, and handling of siblings or shared devices must be documented before recruitment. Participants cannot be masked to AR exposure; where feasible, independent transfer assessors and outcome coders will be masked to allocation.

\subsection{Protocol deviations and safety}
A trial-level record distinguishes assigned condition, delivered condition, technical invalidity, voluntary pause, child-requested stop, and safety override. Safety overrides take precedence over data completeness. Dizziness, discomfort, unsafe movement, or repeated tracking loss trigger the predefined interruption pathway. Such events remain in feasibility and safety denominators even when the associated cognitive trial is uninterpretable.

\subsection{Feasibility and progression criteria}
Before confirmatory recruitment, the team will freeze minimum task comprehension, usable-trial coverage, cross-context timing tolerance, AR tracking stability, acceptable burden, and retention targets. Each threshold requires empirical pilot justification rather than arbitrary retrospective selection. Failure to meet these gates warrants protocol redesign or restriction to a simpler embodied interface, not post hoc exclusion of inconvenient participants.

\subsection{Analysis populations}
The screening population comprises all approached and screened candidates with appropriate permission to retain screening metadata. The randomized population comprises all allocated children. The safety population comprises children exposed to at least one experimental task. The task-quality population comprises prespecified technically interpretable trials and is used only for relevant measurement analyses. The primary transfer analysis retains the randomized denominator and documents missing outcomes separately. These populations must not be conflated when reporting results.